%% =============================================================================
%% PAPER 1 - Section 7: Conclusion
%% =============================================================================

% TODO:
% Summary of findings per RQ:
%   RQ1: Heterogeneous factor structure across strategies
%     — AEN selects different factors for each strategy
%     — PCA confirms that diffuse systematic risk does not span returns
%   RQ2: Alpha persistence across all methodologies
%     — Benchmarks (Duarte, F&H, Active FI), PCA, AEN all yield significant alpha
%     — Robust to: conditional specification (Ferson-Schadt), MPPM, Moreira-Muir,
%       alternative penalized estimators (Adaptive LASSO), SW vs EW weighting
%   RQ3: Evidence for/against slow-moving capital
%     — Duffie scorecard results across subperiods
%     — Common factor interpretable as funding stress
%
% Contribution to literature:
%   — First unified study of 3 textbook FI arbitrages
%   — First academic study of Italian IL bond basis
%   — Disciplined factor selection via AEN + stability selection
%   — Duffie slow-moving capital tested with comprehensive toolkit
%
% Practical implications:
%   — Capital allocation across FI arbitrages
%   — Regulatory implications (balance sheet constraints)
%   — Sustainability of arbitrage under funding stress
%
% Limitations:
%   — Short sample for BTP Italia (~11 years monthly)
%   — Monthly frequency (daily would allow finer analysis)
%   — Small N (3 strategies)
%   — Static betas (potential extension: state-space / Kalman filter)
%
% Future research:
%   — Time-varying betas via state-space models
%   — Extension to other markets / arbitrage types
%   — Higher-frequency analysis